\section{Formal Foundations}\label{sec:foundations}
%==============================================================================

We establish the epistemic foundations of representation theory, then instantiate them to software. The core insight: SSOT is not a software engineering guideline but an epistemic necessity for coherent representation of facts.

\subsection{Epistemic Foundations}\label{sec:epistemic}

Before formalizing codebases and edits, we establish the epistemic structure that underlies all representation systems.

\begin{definition}[Encoding System]\label{def:encoding-system}
An \emph{encoding system} for a fact $F$ is a collection of locations $\{L_1, \ldots, L_n\}$, each capable of holding a value for $F$.
\end{definition}

\begin{definition}[Coherence]\label{def:coherence}
An encoding system is \emph{coherent} iff all locations hold the same value:
\[
\forall i, j: \text{value}(L_i) = \text{value}(L_j)
\]
\end{definition}

\begin{definition}[Incoherence]\label{def:incoherence}
An encoding system is \emph{incoherent} iff some locations disagree:
\[
\exists i, j: \text{value}(L_i) \neq \text{value}(L_j)
\]
\end{definition}

\textbf{The Epistemic Problem.} When an encoding system is incoherent, which value is ``true''? This question has no answer internal to the system.

\begin{theorem}[Oracle Arbitrariness]\label{thm:oracle-arbitrary}
For any incoherent encoding system $S$ and any oracle $O$ that resolves $S$ to a value $v \in S$, there exists a value $v' \in S$ such that $v' \neq v$.
\LH{ORA1}
\end{theorem}

\begin{proof}
By incoherence, $\exists v_1, v_2 \in S: v_1 \neq v_2$. Either $O$ picks $v_1$ (then $v_2$ disagrees) or $O$ doesn't pick $v_1$ (then $v_1$ disagrees).
\end{proof}

\textbf{Interpretation.} This theorem is not about difficulty---it is about \emph{indeterminacy}. Under incoherence, the encoding system does not determine which value is true. Any resolution requires information external to the encodings.

\begin{definition}[Degrees of Freedom]\label{def:dof-epistemic}
The \emph{degrees of freedom} (DOF) of an encoding system is the number of locations that can be modified independently.
\end{definition}

\begin{theorem}[DOF = 1 Guarantees Coherence]\label{thm:dof-one-coherence}
If DOF = 1, then the encoding system is coherent in all reachable states.
\LH{COH1}
\end{theorem}

\begin{proof}
With DOF = 1, exactly one location is independent. All other locations are derived (automatically updated when the source changes). Derived locations cannot diverge from their source. Therefore, all locations hold the value determined by the single independent source. Disagreement is impossible.
\end{proof}

\begin{theorem}[DOF $>$ 1 Permits Incoherence]\label{thm:dof-gt-one-incoherence}
If DOF $> 1$, then incoherent states are reachable.
\LH{COH2}
\end{theorem}

\begin{proof}
With DOF $> 1$, at least two locations are independent. Independent locations can be modified separately. A sequence of edits can set $L_1 = v_1$ and $L_2 = v_2$ where $v_1 \neq v_2$. This is an incoherent state.
\end{proof}

\begin{corollary}[Coherence Forces DOF = 1]\label{cor:coherence-forces-ssot}
If coherence must be guaranteed (no incoherent states reachable), then DOF = 1 is necessary and sufficient.
\LH{COH3}
\end{corollary}

This is the epistemic foundation of SSOT.

\textbf{Information-Theoretic Interpretation.} The DOF = 1 condition has a precise information-theoretic characterization. Rissanen's Minimum Description Length (MDL) principle~\cite{rissanen1978mdl} states that the optimal representation minimizes total description length: $|$model$|$ + $|$data given model$|$. When DOF = 1, the single source \emph{is} the minimal model; all derived locations are the ``data given model'' with zero additional description length (they are fully determined). When DOF $> 1$, the redundant independent locations require explicit synchronization---additional description length that serves no epistemic purpose. Grünwald~\cite{gruenwald2007mdl} provides a comprehensive treatment showing that MDL-optimal representations are unique under mild conditions---precisely what Theorem~\ref{thm:dof-optimal} establishes for DOF.

Heering~\cite{heering2015generative} formalized this connection for software specifically: the \emph{generative complexity} of a program family is the length of the shortest generator. SSOT architectures achieve minimal generative complexity by construction---the single source is the generator, and derived locations are the generated instances. This framing connects our formalization to Kolmogorov complexity theory while remaining constructive (we provide the generator, not just prove its existence).

The following sections instantiate these concepts to software.

%------------------------------------------------------------------------------

\subsection{Software Instantiation: Codebases}\label{sec:edit-space}

\begin{definition}[Codebase]
A \emph{codebase} $C$ is a finite collection of source files, each containing a sequence of syntactic constructs (classes, functions, statements, expressions).
\end{definition}

\begin{definition}[Location]
A \emph{location} $L \in C$ is a syntactically identifiable region of code: a class definition, a function body, a configuration value, a type annotation, etc.
\end{definition}

\begin{definition}[Edit Space]
For a codebase $C$, the \emph{edit space} $E(C)$ is the set of all syntactically valid modifications to $C$. Each edit $\delta \in E(C)$ transforms $C$ into a new codebase $C' = \delta(C)$.
\end{definition}

The edit space is large (exponential in codebase size). But we are not interested in arbitrary edits. We are interested in edits that \emph{change a specific fact}.

\subsection{Facts: Atomic Units of Specification}\label{sec:facts}

\begin{definition}[Fact]\label{def:fact}
A \emph{fact} $F$ is an atomic unit of program specification: a single piece of knowledge that can be independently modified. Facts are the indivisible units of meaning in a specification.
\end{definition}

The granularity of facts is determined by the specification, not the implementation. If two pieces of information must always change together, they constitute a single fact. If they can change independently, they are separate facts.

\noindent\textbf{Examples of facts:}

\begin{center}
\begin{tabular}{lp{8cm}}
\toprule
\textbf{Fact} & \textbf{Description} \\
\midrule
$F_1$: ``threshold = 0.5'' & A configuration value \\
$F_2$: ``\texttt{PNGLoader} handles \texttt{.png}'' & A type-to-handler mapping \\
$F_3$: ``\texttt{validate()} returns \texttt{bool}'' & A method signature \\
$F_4$: ``\texttt{Detector} is a subclass of \texttt{Processor}'' & An inheritance relationship \\
$F_5$: ``\texttt{Config} has field \texttt{name: str}'' & A dataclass field \\
\bottomrule
\end{tabular}
\end{center}

\begin{definition}[Structural Fact]\label{def:structural-fact}
A fact $F$ is \emph{structural} with respect to codebase $C$ iff the locations encoding $F$ are syntactic constituents of type definitions:
\[
\text{structural}(F, C) \Longleftrightarrow \forall L: \text{encodes}(L, F) \rightarrow L \in \text{TypeSyntax}(C)
\]
where $\text{TypeSyntax}(C)$ comprises class declarations, method signatures, inheritance clauses, and attribute definitions.
\end{definition}

\textbf{Key property:} Structural facts are fixed at \emph{definition time}. Once a class is defined, its structure (methods, bases, attributes) cannot change without redefining the class. This is why structural SSOT requires definition-time hooks: the encoding locations are only mutable during class creation.

\textbf{Non-structural facts} (configuration values, runtime state) have encoding locations outside type syntax. They can be modified at runtime without redefining types. SSOT for non-structural facts requires different mechanisms (e.g., reactive bindings, event systems) and is outside the scope of this paper.

\subsection{Encoding: The Correctness Relationship}\label{sec:encoding}

\begin{definition}[Encodes]\label{def:encodes}
Location $L$ \emph{encodes} fact $F$, written $\text{encodes}(L, F)$, iff correctness requires updating $L$ when $F$ changes.

Formally:
\[
\text{encodes}(L, F) \Longleftrightarrow \forall \delta_F: \neg\text{updated}(L, \delta_F) \rightarrow \text{incorrect}(\delta_F(C))
\]

where $\delta_F$ is an edit targeting fact $F$.
\end{definition}

\textbf{Key insight:} This definition is \textbf{forced} by correctness, not chosen. We do not decide what encodes what. Correctness requirements determine it. If failing to update location $L$ when fact $F$ changes produces an incorrect program, then $L$ encodes $F$. This is an objective, observable property.

\begin{example}[Encoding in Practice]\label{ex:encoding}
Consider a type registry:

\begin{verbatim}
# Location L1: Class definition
class PNGLoader(ImageLoader):
    format = "png"

# Location L2: Registry entry
LOADERS = {"png": PNGLoader, "jpg": JPGLoader}

# Location L3: Documentation
# Supported formats: png, jpg
\end{verbatim}

The fact $F$ = ``\texttt{PNGLoader} handles \texttt{png}'' is encoded at:
\begin{itemize}
\tightlist
\item $L_1$: The class definition (primary encoding)
\item $L_2$: The registry dictionary (secondary encoding)
\item $L_3$: The documentation comment (tertiary encoding)
\end{itemize}

If $F$ changes (e.g., to ``\texttt{PNGLoader} handles \texttt{png} and \texttt{apng}''), all three locations must be updated for correctness. The program is incorrect if $L_2$ still says \texttt{\{"png": PNGLoader\}} when the class now handles both formats.
\end{example}

\subsection{Modification Complexity}\label{sec:mod-complexity}

\begin{definition}[Modification Complexity]\label{def:mod-complexity}
\[
M(C, \delta_F) = |\{L \in C : \text{encodes}(L, F)\}|
\]
The number of locations that must be updated when fact $F$ changes.
\end{definition}

Modification complexity is the central metric of this paper. It measures the \emph{cost} of changing a fact. A codebase with $M(C, \delta_F) = 47$ requires 47 edits to correctly implement a change to fact $F$. A codebase with $M(C, \delta_F) = 1$ requires only 1 edit.

\begin{theorem}[Correctness Forcing]\label{thm:correctness-forcing}
$M(C, \delta_F)$ is the \textbf{minimum} number of edits required for correctness. Fewer edits imply an incorrect program.
\LH{BAS1}
\end{theorem}

\begin{proof}
Suppose $M(C, \delta_F) = k$, meaning $k$ locations encode $F$. By Definition~\ref{def:encodes}, each encoding location must be updated when $F$ changes. If only $j < k$ locations are updated, then $k - j$ locations still reflect the old value of $F$. These locations create inconsistencies:

\begin{enumerate}
\tightlist
\item The specification says $F$ has value $v'$ (new)
\item Locations $L_1, \ldots, L_j$ reflect $v'$
\item Locations $L_{j+1}, \ldots, L_k$ reflect $v$ (old)
\end{enumerate}

By Definition~\ref{def:encodes}, the program is incorrect. Therefore, all $k$ locations must be updated, and $k$ is the minimum.
\end{proof}

\subsection{Independence and Degrees of Freedom}\label{sec:dof}

Not all encoding locations are created equal. Some are \emph{derived} from others.

\begin{definition}[Independent Locations]\label{def:independent}
Locations $L_1, L_2$ are \emph{independent} for fact $F$ iff they can diverge. Updating $L_1$ does not automatically update $L_2$, and vice versa.

Formally: $L_1$ and $L_2$ are independent iff there exists a sequence of edits that makes $L_1$ and $L_2$ encode different values for $F$.
\end{definition}

\begin{definition}[Derived Location]\label{def:derived}
Location $L_{\text{derived}}$ is \emph{derived from} $L_{\text{source}}$ iff updating $L_{\text{source}}$ automatically updates $L_{\text{derived}}$. Derived locations are not independent of their sources.
\end{definition}

\begin{example}[Independent vs. Derived]\label{ex:independence}
Consider two architectures for the type registry:

\textbf{Architecture A (independent locations):}
\begin{verbatim}
# L1: Class definition
class PNGLoader(ImageLoader): ...

# L2: Manual registry (independent of L1)
LOADERS = {"png": PNGLoader}
\end{verbatim}

Here $L_1$ and $L_2$ are independent. A developer can change $L_1$ without updating $L_2$, causing inconsistency.

\textbf{Architecture B (derived location):}
\begin{verbatim}
# L1: Class definition with registration
class PNGLoader(ImageLoader):
    format = "png"

# L2: Derived registry (computed from L1)
LOADERS = {cls.format: cls for cls in ImageLoader.__subclasses__()}
\end{verbatim}

Here $L_2$ is derived from $L_1$. Updating the class definition automatically updates the registry. They cannot diverge.
\end{example}

\begin{definition}[Degrees of Freedom]\label{def:dof}
\[
\text{DOF}(C, F) = |\{L \in C : \text{encodes}(L, F) \land \text{independent}(L)\}|
\]
The number of \emph{independent} locations encoding fact $F$.
\end{definition}

DOF is the key metric. Modification complexity $M$ counts all encoding locations. DOF counts only the independent ones. If all but one encoding location is derived, DOF = 1 even though $M$ may be large.

\begin{theorem}[DOF = Incoherence Potential]\label{thm:dof-inconsistency}
$\text{DOF}(C, F) = k$ implies $k$ different values for $F$ can coexist in $C$ simultaneously. With $k > 1$, incoherent states are reachable.
\LH{BAS2}
\end{theorem}

\begin{proof}
Each independent location can hold a different value. By Definition~\ref{def:independent}, no constraint forces agreement between independent locations. Therefore, $k$ independent locations can hold $k$ distinct values. This is an instance of Theorem~\ref{thm:dof-gt-one-incoherence} applied to software.
\end{proof}

\begin{corollary}[DOF $>$ 1 Implies Incoherence Risk]\label{cor:dof-risk}
$\text{DOF}(C, F) > 1$ implies incoherent states are reachable. The codebase can enter a state where different locations encode different values for the same fact.
\LH{BAS3}
\end{corollary}

\subsection{The DOF Lattice}\label{sec:dof-lattice}

DOF values form a lattice with distinct epistemic meanings:

\begin{center}
\begin{tabular}{cl}
\toprule
\textbf{DOF} & \textbf{Epistemic Status} \\
\midrule
0 & Fact $F$ is not represented (no encoding) \\
1 & Coherence guaranteed (truth determinate) \\
$k > 1$ & Incoherence possible (truth indeterminate under divergence) \\
\bottomrule
\end{tabular}
\end{center}

\begin{theorem}[DOF = 1 is Uniquely Coherent]\label{thm:dof-optimal}
For any fact $F$ that must be encoded, $\text{DOF}(C, F) = 1$ is the unique value guaranteeing coherence:
\begin{enumerate}
\tightlist
\item DOF = 0: Fact is not represented
\item DOF = 1: Coherence guaranteed (by Theorem~\ref{thm:dof-one-coherence})
\item DOF $>$ 1: Incoherence reachable (by Theorem~\ref{thm:dof-gt-one-incoherence})
\end{enumerate}
\LH{INS1}
\end{theorem}

\begin{proof}
This is a direct instantiation of Corollary~\ref{cor:coherence-forces-ssot} to software:
\begin{enumerate}
\item DOF = 0 means no location encodes $F$. The fact is unrepresented.
\item DOF = 1 means exactly one independent location. All other encodings are derived. Divergence is impossible. Coherence is guaranteed.
\item DOF $>$ 1 means multiple independent locations. By Corollary~\ref{cor:dof-risk}, they can diverge. Incoherence is reachable.
\end{enumerate}

Only DOF = 1 achieves coherent representation. This is epistemic necessity, not design preference.
\end{proof}

%==============================================================================
